%2multibyte Version: 5.50.0.2953 CodePage: 1252

\documentclass{article}
%%%%%%%%%%%%%%%%%%%%%%%%%%%%%%%%%%%%%%%%%%%%%%%%%%%%%%%%%%%%%%%%%%%%%%%%%%%%%%%%%%%%%%%%%%%%%%%%%%%%%%%%%%%%%%%%%%%%%%%%%%%%%%%%%%%%%%%%%%%%%%%%%%%%%%%%%%%%%%%%%%%%%%%%%%%%%%%%%%%%%%%%%%%%%%%%%%%%%%%%%%%%%%%%%%%%%%%%%%%%%%%%%%%%%%%%%%%%%%%%%%%%%%%%%%%%
\usepackage{amsfonts}
\usepackage{amsmath}

\setcounter{MaxMatrixCols}{10}
%TCIDATA{OutputFilter=LATEX.DLL}
%TCIDATA{Version=5.50.0.2953}
%TCIDATA{Codepage=1252}
%TCIDATA{<META NAME="SaveForMode" CONTENT="1">}
%TCIDATA{BibliographyScheme=Manual}
%TCIDATA{Created=Tuesday, September 13, 2016 11:26:35}
%TCIDATA{LastRevised=Monday, December 09, 2019 17:17:30}
%TCIDATA{<META NAME="GraphicsSave" CONTENT="32">}
%TCIDATA{<META NAME="DocumentShell" CONTENT="Standard LaTeX\Blank - Standard LaTeX Article">}
%TCIDATA{Language=American English}
%TCIDATA{CSTFile=40 LaTeX article.cst}

\newtheorem{theorem}{Theorem}
\newtheorem{acknowledgement}[theorem]{Acknowledgement}
\newtheorem{algorithm}[theorem]{Algorithm}
\newtheorem{axiom}[theorem]{Axiom}
\newtheorem{case}[theorem]{Case}
\newtheorem{claim}[theorem]{Claim}
\newtheorem{conclusion}[theorem]{Conclusion}
\newtheorem{condition}[theorem]{Condition}
\newtheorem{conjecture}[theorem]{Conjecture}
\newtheorem{corollary}[theorem]{Corollary}
\newtheorem{criterion}[theorem]{Criterion}
\newtheorem{definition}[theorem]{Definition}
\newtheorem{example}[theorem]{Example}
\newtheorem{exercise}[theorem]{Exercise}
\newtheorem{lemma}[theorem]{Lemma}
\newtheorem{notation}[theorem]{Notation}
\newtheorem{problem}[theorem]{Problem}
\newtheorem{proposition}[theorem]{Proposition}
\newtheorem{remark}[theorem]{Remark}
\newtheorem{solution}[theorem]{Solution}
\newtheorem{summary}[theorem]{Summary}
\newenvironment{proof}[1][Proof]{\noindent\textbf{#1.} }{\ \rule{0.5em}{0.5em}}
\input{tcilatex}
\begin{document}


\begin{center}
{\Large Final Exam}
\end{center}

\textbf{Question 1}

(a) Explain what the difference is between an unbiased and a consistent
estimator. Explain why these are both desirable properties for estimators to
have.

(b) Define the size and the power of a hypothesis test and discuss how these
relate to Type I and II errors in testing.

(c) Define the score vector and Fisher information matrix and state the Cram%
\'{e}r Rao theorem.

\bigskip

\textbf{Question 2}

Consider the linear regression model $y=X\beta +\varepsilon ,$ where $y$ is
an $n\times 1$ vector of dependent variables, $X$ is an $n\times k$ matrix
of regressors of rank $k$, $\beta $ is a $k\times 1$ vector of unknown
parameters, and $\varepsilon $ is an $n\times 1$ vector of unobserved random
variables. The standard assumptions provided in the lecture are assumed to
hold, particularly $\mathbb{E}\left[ \varepsilon |X\right] =0$ and $\mathbb{E%
}\left[ \varepsilon \varepsilon ^{\prime }|X\right] =\sigma ^{2}I_{n}.$

(a) Suppose there are two explanatory variables, $k=2,$ so that $%
y=X_{1}\beta _{1}+X_{2}\beta _{2}+\varepsilon .$ State and prove the
Frisch-Waugh-Lovell theorem for the OLS estimator $\hat{\beta}_{1}.$

(b) Consider two economists estimating the following regression models:
Economist 1 estimates $y=X_{1}\beta _{1}+\varepsilon $ while Economist 2
also includes $X_{2}$ and estimates instead $y=X_{1}\beta _{1}+X_{2}\beta
_{2}+\varepsilon .$ What are the estimators of the two economists for $\beta
_{1}$?

(c) Given that the true data generating process is $y=X_{1}\beta
_{1}+X_{2}\beta _{2}+\varepsilon ,$ show that that the estimator for $\beta
_{1}$ of Economist 2 is unbiased.

(d) Show that the estimator for $\beta _{1}$ of Economist 1 is biased.
Discuss this result in relation to omitted variable bias.

(e)\ Now suppose the true data generating process is $y=X_{1}\beta
_{1}+\varepsilon ,$ so $\beta _{2}=0.$ Are the two estimators of the two
economists unbiased?

(f) Given that the true $\beta _{2}=0,$ compare the variance of $\hat{\beta}%
_{1}$ for the two economists. What can you conclude about the tradeoff
between bias and variance when deciding which explanatory variables to
include?

\bigskip

\textbf{Question 3 }

Consider the Gaussian linear model%
\begin{equation*}
y=X\beta +\varepsilon 
\end{equation*}%
where $\beta \in \mathbb{R}^{K}$, $X$ is an $n\times k$ matrix of regressors
of rank $k,$ and $\varepsilon |X\sim \mathcal{N}\left( 0,\sigma ^{2}\Omega
\right) $ for some known positive definite symmetric matrix $\Omega $.

(a) Write down the log-likelihood function of the sample $y$, recalling that
the density of an $m\times 1$ vector $z\sim \mathcal{N}\left( \mu ,\Sigma
\right) $ is given by 
\begin{equation*}
f\left( z;\mu ,\Sigma \right) =\left( 2\pi \right) ^{-m/2}\det \left( \Sigma
\right) ^{-1/2}\exp \left\{ -\frac{1}{2}\left( z-\mu \right) ^{\prime
}\Sigma ^{-1}\left( z-\mu \right) \right\} .
\end{equation*}

(b) Show that the maximum likelihood estimator (MLE) of $\beta $ coincides
with the generalised least squares estimator $\hat{\beta}^{GLS}$ of $\beta $.

(c) Show that the MLE of $\beta $ is unbiased and calculate its covariance
matrix conditional on $X$.

(d) Find the maximum likelihood estimator (MLE) of $\sigma ^{2}$.

(e) Compute the Fisher Information matrix conditional on $X$. 

(f) Show that the variance of $\hat{\beta}^{MLE}$ conditional on $X$
achieves the Cram\'{e}r Rao Lower Bound. Is this result the same as the
Gauss Markov theorem for GLS?

\end{document}
